% Vietnamese translation for OI Public Library
% Source-File: IOI中国国家候选队论文/2003/周源/周源.doc
\documentclass[11pt]{article}
\usepackage{oipl-vn}

\title{Ứng dụng tư tưởng biểu diễn nhỏ nhất trong bài toán đẳng cấu vòng của chuỗi}
\author{Zhou Yuan}
\date{2003}

\begin{document}
\maketitle

\origtitle{Application of minimal representation to cyclic string isomorphism}
\sourcepath{IOI China candidate team papers / 2003 / Zhou Yuan / Zhou Yuan.doc}

\begin{abstract}
Biểu diễn nhỏ nhất có nhiều ứng dụng quan trọng trong loại bỏ trạng thái trùng
lặp khi tìm kiếm, kiểm tra đẳng cấu đồ thị và nhiều bài toán khác. Bài viết
lấy bài toán kiểm tra hai chuỗi có đẳng cấu vòng hay không làm trung tâm. Sau
khi chỉ ra lời giải khớp mẫu tuyến tính quen thuộc, bài viết đưa vào tư tưởng
``biểu diễn nhỏ nhất'', định nghĩa nó một cách hệ thống, rồi dùng tư tưởng đó
xây dựng một thuật toán tự nhiên và gọn hơn cho bài toán.
\end{abstract}

\noindent\textbf{Từ khóa:} chuỗi, đẳng cấu vòng, khớp mẫu, biểu diễn nhỏ
nhất.

\section{Dẫn nhập bài toán}

Vì bài viết chủ yếu bàn về thuật toán trên chuỗi, ta dùng các ký hiệu sau.

\begin{itemize}
  \item \(|s|\) là độ dài chuỗi \(s\).
  \item \(s_i\) là ký tự thứ \(i\) của \(s\), với \(1\le i\le |s|\).
  \item \(s[i..j]\) là chuỗi con từ ký tự thứ \(i\) đến ký tự thứ \(j\),
  với \(1\le i\le j\le |s|\). Quy ước \(s[i..i]=s_i\).
  \item Một lần quay vòng của \(s\) đưa ký tự đầu xuống cuối:
  \[
    \operatorname{rot}(s)=s[2..|s|]s_1.
  \]
  Lần quay thứ \(k\) được ký hiệu là \(\operatorname{rot}^k(s)\), và
  \(\operatorname{rot}^0(s)=s\).
  \item Nếu tồn tại \(k\) sao cho \(t=\operatorname{rot}^k(s)\), ta nói
  \(s\) và \(t\) là \term{đẳng cấu vòng}.
  \item Với hai ánh xạ \(f,g\), ký hiệu \(g\circ f\) là phép hợp thành:
  \((g\circ f)(x)=g(f(x))\). Ký hiệu này dùng cho phần định nghĩa trừu tượng.
\end{itemize}

Bài toán cần giải là: cho hai chuỗi \(s_1,s_2\) có cùng độ dài
\[
  N=|s_1|=|s_2|,
\]
hãy xác định chúng có đẳng cấu vòng hay không.

\section{Thuật toán liệt kê và thuật toán khớp mẫu}

\subsection{Liệt kê trực tiếp}

Một chuỗi \(s_1\) có nhiều nhất \(N\) chuỗi quay vòng khác nhau:
\[
  s_1,\operatorname{rot}(s_1),\ldots,\operatorname{rot}^{N-1}(s_1).
\]
Vì vậy có thể liệt kê từng chuỗi quay vòng rồi so sánh với \(s_2\). Cách này
dễ nghĩ và dễ cài đặt, nhưng thời gian là \(O(N^2)\). Khi \(N\) lên tới vài
trăm nghìn hoặc vài triệu, nó không còn phù hợp.

\subsection{Chuyển thành khớp mẫu}

Quan sát quá trình liệt kê, ta thấy bản chất của nó là tìm xem \(s_2\) có
xuất hiện trong một chuỗi vòng sinh từ \(s_1\) hay không. Nếu loại bỏ yếu tố
``vòng'', bài toán trở thành khớp mẫu thông thường.

Cách loại bỏ rất đơn giản: ghép \(s_1\) với chính nó:
\[
  S=s_1s_1.
\]
Khi đó mọi chuỗi quay vòng của \(s_1\) đều xuất hiện như một chuỗi con độ dài
\(N\) trong \(S\). Vì vậy \(s_1\) và \(s_2\) đẳng cấu vòng khi và chỉ khi
\(s_2\) xuất hiện trong \(s_1s_1\). Dùng KMP hoặc một thuật toán khớp mẫu tuyến
tính khác, ta giải được bài toán trong \(O(N)\) thời gian.

\subsection{Nhận xét}

Thuật toán khớp mẫu đã đạt độ phức tạp tốt, nhưng nó vẫn có nhược điểm thực
tế: KMP nổi tiếng là khó nhớ, đặc biệt là phần tính mảng \(\mathit{next}\),
và điều đó làm giảm khả năng mở rộng ý tưởng. Vì vậy vẫn đáng tìm một cách
nghĩ khác, đơn giản hơn về mặt cấu trúc.

\section{Tư tưởng biểu diễn nhỏ nhất}

\subsection{Gợi ý từ một ví dụ}

Xét bài toán rất đơn giản: cho hai dãy số, không xét thứ tự, hãy kiểm tra
chúng có giống nhau hay không. Nếu liệt kê mọi hoán vị thì có tới \(N!\) dạng
biểu diễn khác nhau, rõ ràng không hợp lý. Cách tự nhiên hơn là sắp xếp cả hai
dãy rồi so sánh kết quả.

Ví dụ nhỏ này cho ta một tư tưởng quan trọng: khi một đối tượng có nhiều dạng
biểu diễn khác nhau, và ta muốn biết hai đối tượng có thể biến đổi về cùng
một dạng hay không, hãy đưa mỗi đối tượng về dạng nhỏ nhất theo một thứ tự cố
định, rồi chỉ cần so sánh hai dạng nhỏ nhất đó.

\subsection{Định nghĩa trừu tượng}

Gọi \(V\) là tập các sự vật, \(F\) là tập các phép biến đổi trên \(V\). Mỗi
\(f\in F\) là một ánh xạ
\[
  f:V\to V.
\]
Nếu với hai phần tử \(a,b\in V\), tồn tại một dãy phép biến đổi trong \(F\)
sao cho hợp thành của chúng đưa \(a\) thành \(b\), ta nói \(a\) và \(b\)
\term{giống nhau về bản chất} theo \(F\).

Tập phép biến đổi \(F\) được giả sử có hai tính chất:

\begin{enumerate}
  \item Với mọi \(a\in V\), áp dụng một dãy phép biến đổi trong \(F\) vẫn cho
  một phần tử thuộc \(V\).
  \item Nếu có \(f\in F\) sao cho \(f(a)=b\), thì tồn tại một hoặc nhiều phép
  biến đổi trong \(F\) mà hợp thành của chúng đưa \(b\) trở lại \(a\).
\end{enumerate}

Hai tính chất này làm cho quan hệ ``giống nhau về bản chất'' có tính phản xạ,
đối xứng và bắc cầu. Nói cách khác, nó chia \(V\) thành các lớp tương đương.

Tư tưởng biểu diễn nhỏ nhất áp dụng như sau. Trước hết quy định một thứ tự
trên \(V\). Với mỗi phần tử \(a\), xét tất cả các phần tử có thể thu được từ
\(a\) bằng các phép biến đổi trong \(F\), rồi chọn phần tử nhỏ nhất trong số
đó, ký hiệu \(\min_F(a)\). Khi đó
\[
  a \text{ và } b \text{ giống nhau về bản chất}
  \quad\Longleftrightarrow\quad
  \min_F(a)=\min_F(b).
\]

\section{Áp dụng cho chuỗi đẳng cấu vòng}

Trong bài toán hiện tại, \(V\) là tập các chuỗi, \(F\) là tập các phép quay
vòng, còn thứ tự là thứ tự từ điển giữa các chuỗi. Biểu diễn nhỏ nhất của một
chuỗi là chuỗi nhỏ nhất theo thứ tự từ điển trong tất cả các phép quay vòng
của nó.

Cách trực tiếp là tìm biểu diễn nhỏ nhất của \(s_1\), tìm biểu diễn nhỏ nhất
của \(s_2\), rồi so sánh. Tuy nhiên bài toán tìm biểu diễn nhỏ nhất của một
chuỗi tuy có thuật toán tuyến tính, nhưng cách trình bày thường không đơn
giản hơn KMP. Nếu làm máy móc như vậy, ta chưa đạt mục tiêu tìm một thuật
toán dễ hiểu hơn.

Ta dùng tư tưởng biểu diễn nhỏ nhất theo cách trực tiếp hơn. Đặt
\[
  u=s_1s_1,\qquad w=s_2s_2,
\]
và dùng hai con trỏ \(i,j\), ban đầu đều bằng \(1\), lần lượt chỉ vào vị trí
khởi đầu của một phép quay trong \(u\) và \(w\). Gọi
\[
  p(s)
\]
là vị trí nhỏ nhất trong chuỗi gốc mà từ đó bắt đầu một biểu diễn nhỏ nhất
của \(s\). Nếu \(s_1\) và \(s_2\) đẳng cấu vòng, thì mục tiêu ngầm của ta là
đưa \(i\) và \(j\) tới \(p(s_1)\) và \(p(s_2)\).

Giả sử từ \(i\) và \(j\), ta so sánh lần lượt:
\[
  u_{i}=w_{j},\ u_{i+1}=w_{j+1},\ldots,
  u_{i+k-1}=w_{j+k-1},
\]
nhưng đến ký tự kế tiếp thì khác:
\[
  u_{i+k}\ne w_{j+k}.
\]

Nếu \(u_{i+k}>w_{j+k}\), mọi phép quay của \(s_1\) bắt đầu tại các vị trí
\[
  i,i+1,\ldots,i+k
\]
đều không thể là biểu diễn nhỏ nhất. Lý do là chúng cùng có một đoạn đầu khớp
với phần tương ứng của \(w\), nhưng tại ký tự đầu tiên khác nhau thì phía
\(u\) lớn hơn. Do đó có thể bỏ qua cả đoạn ấy và đặt
\[
  i\leftarrow i+k+1.
\]
Ngược lại, nếu \(u_{i+k}<w_{j+k}\), ta bỏ qua đoạn tương ứng của \(w\):
\[
  j\leftarrow j+k+1.
\]

Nếu không tìm được vị trí khác nhau trong \(N\) ký tự liên tiếp, tức là
\[
  u[i..i+N-1]=w[j..j+N-1],
\]
thì hai chuỗi đẳng cấu vòng.

\subsection{Thuật toán}

Thuật toán có thể mô tả như sau.

\begin{enumerate}
  \item Đặt \(u=s_1s_1\), \(w=s_2s_2\), \(i=j=1\).
  \item Nếu \(i>N\) hoặc \(j>N\), trả lời rằng hai chuỗi không đẳng cấu vòng.
  \item Tìm \(k\) nhỏ nhất sao cho \(0\le k<N\) và
  \(u_{i+k}\ne w_{j+k}\).
  \item Nếu không tồn tại \(k\), trả lời rằng hai chuỗi đẳng cấu vòng.
  \item Nếu \(u_{i+k}>w_{j+k}\), đặt \(i\leftarrow i+k+1\); nếu ngược lại,
  đặt \(j\leftarrow j+k+1\). Quay lại bước 2.
\end{enumerate}

Mỗi lần thất bại, ít nhất một con trỏ tăng thêm \(k+1\), còn các vị trí bị bỏ
qua sẽ không bao giờ được xét lại như điểm bắt đầu. Vì \(i\) và \(j\) đều
không vượt quá \(N+1\), tổng số lần so sánh là tuyến tính. Thuật toán dùng
\(O(N)\) thời gian và \(O(N)\) bộ nhớ nếu thật sự tạo \(u,w\); có thể giảm bộ
nhớ xuống \(O(1)\) bằng cách truy cập chỉ số theo modulo \(N\).

\subsection{Mô phỏng}

Xét hai chuỗi đẳng cấu vòng. Khi thuật toán chạy, mỗi vòng lặp so sánh hai
phép quay hiện tại. Nếu gặp ký tự khác nhau, bên có ký tự lớn hơn tại vị trí
đầu tiên khác nhau chắc chắn không thể chứa biểu diễn nhỏ nhất trong toàn bộ
đoạn vừa so sánh, nên con trỏ bên đó nhảy qua đoạn. Sau vài lần nhảy, hai con
trỏ sẽ cùng trỏ tới hai phép quay có cùng biểu diễn nhỏ nhất; lúc đó \(N\) ký
tự liên tiếp khớp nhau và thuật toán trả lời đúng.

Điểm quan trọng ở đây là ta không cần trực tiếp dựng biểu diễn nhỏ nhất của
từng chuỗi. Ta chỉ dùng sự tồn tại của nó để biện minh cho các bước loại bỏ
hàng loạt vị trí khởi đầu.

\section{Tổng kết}

Sau quá trình phân tích, ta thu được một thuật toán tuyến tính khác hẳn bản
chất với khớp mẫu. Trong bài toán này, tư tưởng biểu diễn nhỏ nhất giúp ta
nhìn bài toán từ phía ``dạng chuẩn'' của mỗi lớp tương đương, rồi từ đó xây
dựng một quy tắc nhảy con trỏ đơn giản.

Biểu diễn nhỏ nhất là một tư tưởng thường gặp khi cần kiểm tra hai đối tượng
có cùng bản chất dưới một nhóm phép biến đổi hay không. Nó xuất hiện trong
loại bỏ trạng thái trùng lặp khi tìm kiếm, trong kiểm tra đẳng cấu đồ thị và
trong nhiều bài toán mô hình hóa khác. Tinh hoa của tư tưởng này là đưa
khái niệm ``thứ tự'' vào tập các biểu diễn khác nhau, chọn một đại diện chuẩn
nhỏ nhất hoặc lớn nhất, rồi so sánh các đại diện chuẩn đó.

Dù vậy, trong thi tin học hiện đại, bài toán hiếm khi cho phép áp dụng máy
móc một thuật toán kinh điển. Điều đáng học ở đây không chỉ là bản thân thuật
toán cho chuỗi vòng, mà là cách dùng tư tưởng biểu diễn nhỏ nhất như một công
cụ phân tích. Khi hiểu sâu cấu trúc của bài toán và dám thay đổi góc nhìn, ta
có thể biến quan hệ rối rắm thành có trật tự, và biến lời giải phức tạp thành
một quy trình gọn gàng.

\end{document}
